\documentclass[11pt,a4paper]{article}
\usepackage[T1]{fontenc}
\usepackage[utf8]{inputenc}
\usepackage{lmodern}
\usepackage[margin=25mm]{geometry}
\usepackage{amsmath,amssymb,booktabs,tabularx,xurl,xcolor,hyperref}
\hypersetup{colorlinks=true,linkcolor=blue!50!black,urlcolor=blue!50!black,
 pdftitle={Typed Scalar Extraction and Observation Boundaries},pdfauthor={Mettapedia Working Notes}}
\setlength{\emergencystretch}{3em}
\setlength{\parindent}{0pt}
\setlength{\parskip}{5pt}
\newcommand{\code}[1]{\nolinkurl{#1}}
\title{Typed Scalar Extraction and Observation Boundaries}
\author{Mettapedia Working Notes}
\date{1 November 2026}
\begin{document}
\maketitle
\begin{abstract}
A proof-producing extractor translates a fragment of existing PeTTa atoms
into a result-indexed scalar representation. It handles signed integer and
symbol data, Boolean comparisons, conditional branches and flat tuple
results. Successful evaluation refines the existing source evaluator,
including singleton answer multiplicity, unchanged state and empty I/O.
Callback preparation uses the actual source clause selector and captures,
while certificate revalidation checks local program dependencies. A runtime
counterexample illustrates the central information boundary: Boolean
comparisons encoded as integer zero or one are adequate for branch selection
but not for ordinary value observation. These are source and representation
laws, not a verification of the native C recognizer or the whole fold builtin.
\end{abstract}

\section{A result encoding must name its observers}

The native scalar-fold audit found a comparison node producing a register
marked as an integer, with value zero or one. The recognizer also admitted
comparison nodes as ordinary returned expressions, tuple fields and operands
of other scalar operations. On the inspected binary, native and reference
routes consequently disagreed:

\begin{center}
\begin{tabularx}{\linewidth}{@{}>{\raggedright\arraybackslash}X>{\raggedright\arraybackslash}X>{\raggedright\arraybackslash}X@{}}
\toprule
Observation & Ordinary result & Native scalar result\\
\midrule
Returned comparison & Boolean true or false & Integer 1 or 0\\
Comparison in a tuple & \texttt{(tag true)} & \texttt{(tag 1)}\\
Comparison consumed by addition & Caught arithmetic type error & Successful integer 1\\
Equality of a true comparison with integer 1 & False & Integer 1\\
Comparison used only to choose 7 or 9 & 7 & 7\\
\bottomrule
\end{tabularx}
\end{center}

The positive branch control is important. Let $e(b)=1$ when $b$ is true and
$e(b)=0$ otherwise. Then branching on $e(b)\ne0$ preserves the selected branch.
It does not follow that publishing $e(b)$ preserves the source value. The
formalization proves that no function $d:\mathbb Z\to\mathsf{Atom}$ can satisfy
both
\[
 \forall n,\ d(n)=\mathsf{integer}(n)
 \quad\text{and}\quad
 \forall b,\ d(e(b))=\mathsf{boolean}(b).
\]
The two required meanings of $d(1)$ conflict. A static result sort or runtime
tag must retain the distinction. This is an exact information-loss boundary,
not a claim about a globally minimum encoding.

\section{The typed scalar core}

\code{Machines.TypedScalar} separates two result kinds:
\[
 \mathsf{Value}(\mathsf{data})=\mathbb Z+\mathsf{String},
 \qquad
 \mathsf{Value}(\mathsf{test})=\mathsf{Bool}.
\]
Register reads return data. Their integer/symbol alternative may vary between
executions; the kind system does not require a permanently integer register.
Arithmetic requires data operands and dynamically refuses non-integers.
Comparisons return tests. Equality can compare either kind and keeps Boolean
versus integer values distinct. A conditional requires a test and evaluates
only its selected branch.

The core interpreter is parameterized by its register environment. The read
footprint theorem proves that agreement on all referenced registers suffices
for equal results. This is a conservative static footprint, including both
conditional branches, not a minimum dynamic read set. Repeated reads retain
one shared environment and therefore cannot acquire independent values.

The result index supplies information that can be erased in some physical
positions and must be restored in others. A branch selector can use a Boolean
bit. A returned field must be encoded according to its kind. The typed core
does not silently allow a predicate to feed integer addition.

\section{Extraction from existing source atoms}

\code{PeTTa.Bridges.ScalarSource} consumes the existing \code{OSLFCore.Atom}
carrier. Its supported forms include signed integer literals, inert symbol
literals, Boolean literals, captured data variables, addition, subtraction,
integer orders, equality and conditional expressions. The extractor returns
the inferred result kind, an expression at that kind and a structural
elaboration derivation.

The correspondence theorem is proved by induction on that derivation, using
the existing PeTTa primitive, argument-continuation and conditional laws. If
the extracted expression succeeds with value $v$, the original expression
has a \code{PureReturns} derivation for the correctly encoded $v$, with
unchanged state and exactly one answer. This is not established by defining
the source evaluator to be the extracted interpreter.

The completed-run theorem then characterizes every completed source run:
its state is unchanged, its answer list is the singleton encoded value, and
its input and output observations are empty. It also establishes existence
of sufficient source fuel. It does not supply equality between native work
and evaluator transitions, or permission to bypass a source-visible budget.

Unknown variables, effects, nondeterministic answer constructors and invalid
sort combinations decline extraction. Decline is not a source-language error:
the ordinary evaluator remains responsible for that expression. In particular,
Boolean arithmetic declines rather than being silently coerced to integer
arithmetic. Both branches are checked during extraction, even though only
the selected branch is evaluated.

The source model is the repository's explicit PeTTa operational profile.
Primitive heads have its native meaning. The extractor additionally requires
that no equation redefine an admitted primitive head. Inert-symbol evidence
excludes source-callable heads and control forms. Correspondence with all
native profiles, foreign primitive registrations and their revision counters
remains an implementation bridge.

\section{Flat tuples and live callability}

\code{ScalarBody} compiles scalar results, flat tuples and result-level
conditional branches. Tuple fields carry their individual result kinds, so
a comparison field materializes a Boolean rather than an integer. The
source proof composes the existing left-to-right data-argument continuation;
it preserves field order and does not quotient answer occurrences.

Static and live head admission are distinct. A source tuple may start with
a variable, an integer or an inert symbol. After its fields are computed, the
body evaluator checks the actual head against the captured program. A symbol
that now names a callable function declines. The negative control keeps the
same variable-headed tuple syntax but adds a definition for its produced
head; static body compilation succeeds while live evaluation declines.

This boundary is relevant to caching and lifetime management. The source
certificate records its program and capture assumptions, while the produced
head still has a live obligation. Private field evaluation cannot justify
publishing a partially accepted tuple.

\section{Call-specific selection and captured registers}

\code{ScalarInvocation.prepare} uses the existing \code{Eval.clauses}
selector on the supplied accumulator and item values. It requires exactly
one selected evaluation body. It then extracts scalar registers from the
selected substitution and compiles that body. Registers are justified by the
actual lookup and scalar-decoding results, including bindings obtained from
tuple patterns.

The main theorem concerns an actual two-variable source call whose arguments
are already captured values. An accepted prepared result is that call's
singleton answer with unchanged state and no I/O. Primitive/control-head
exclusions, registration and source clause selection remain explicit checked
conditions.

This permits a useful widening relative to a globally single-equation rule.
An additional equation whose argument pattern cannot match the current call
does not disable extraction. Two matching answer-producing equations decline.
The controls also distinguish a repeated-variable tuple pattern applied to
equal versus unequal arguments. These cases use the existing matcher rather
than treating tuple holes as independently scoped variables.

No rule here removes an arbitrary applicable equation merely because it is
expected to be unproductive. Proving a selected body's empty answer stream
would be an additional analysis. Full fold invocation, changing accumulator
shapes across iterations and non-single-answer callbacks also require the
appropriate surrounding protocol.

\section{Local certificate revalidation}

\code{ScalarDependencies} computes the primitive and inert-symbol head
requirements of an extracted scalar expression. A transport theorem
re-establishes its source elaboration under a new program satisfying those
requirements. The executable revalidator returns a certificate indexed by
that new program, rather than attaching an old proof to a new revision.

Consequently an unrelated equation need not invalidate a scalar certificate,
whereas redefining its comparison primitive or activating a required inert
symbol does. These are semantic dependency tests. A cheap revision-stamp
implementation needs a separate proof that its stamps cover the same
dependencies. The scalar cache law also does not erase the live tuple-head
test or the callback's selected-clause obligation.

The resulting region judgment combines several kinds of requisite information:
\begin{center}
\begin{tabularx}{\linewidth}{@{}lX@{}}
\toprule
Evidence & What it preserves\\
\midrule
Result kind & Boolean versus integer/symbol interpretation\\
Captured environment & Shared variable values and tuple-pattern bindings\\
Program dependencies & The head meanings used by the certificate\\
Live head check & Data versus callable result interpretation\\
Completion observation & Singleton result, state and I/O behavior\\
\bottomrule
\end{tabularx}
\end{center}
This is a concrete connection between dependent result types and optimization
admission. It does not select a universal Prime calculus or identify every
future observer with this source profile.

\section{Qualification and remaining bridges}

\code{ScalarExtractionAudit} imports the generic typed core and six source
modules, prints the principal axiom dependencies and emits declaration-level
foundation manifests. Positive and negative controls cover scalar and tuple
Booleans, nested comparisons, symbol arithmetic, effects, repeated captures,
callability changes, irrelevant and duplicate clauses, and local revalidation.
The seven mathematical modules pass strict warning-as-error checks. Their
manifests cover 585 checked declaration bodies, with eight nonlogical
recursion companions listed separately, no admissions and no added axioms.
Transitive primitive dependencies are limited to \code{propext},
\code{Classical.choice} and \code{Quot.sound}.

The native discrepancy has a separate differential fixture and captured
outputs. A successful Lean audit does not repair that C path. Native Boolean
tagging or predicate-only admission needs its own regression and mutation
qualification.

The next correspondence layer is a bounded, typed register implementation
with checked signed arithmetic, correct result materialization and refusal
before public effects. The existing private-fold law supplies the restart
boundary, but complete builtin invocation, source budget and yield accounting,
allocation failure and native lifetime obligations are not discharged by
the callback extraction theorem.
\end{document}
